\exercisesection{8}

\begin{enumerate}
\def\labelenumi{\arabic{enumi}.}
\item
  设 K 为域，A = K{[}{[}X, Y, Z{]}{]} 为 K 上三个不定元的形式幂级数环；令 \(v'\)（分别 \(v''\)）为 A 上取值于按字典序排序的群 \(Z \times Z\) 的赋值，满足 \(v'(X) = (1, 0)\)、\(v'(Y) = (0, 1)\)，且 \(v'\) 在 K{[}{[}Z{]}{]} 上是非真赋值（分别地，\(v''(X) = (1, 0)\)、\(v''\) 在 K{[}{[}Y{]}{]} 上是非真赋值、\(v''(Z) = (0, 1)\)）。令 \(\sigma\) 为保持 K 与 X 不变且满足 \(\sigma(Y) = Z\)、\(\sigma(Z) = Y\) 的 A 的自同构；若 B 是 A 中由 \(\sigma\) 不变元素组成的子环，E（分别 F）为 A（分别 B）的分式域，则 \(v'\) 与 \(v''\)（典范延拓到 E 后）在 F 上的限制相同，F 关于 v 定义的拓扑完备，而 E 是 F 的二次扩张；E 上的两个赋值 \(v'\)、\(v''\) 相依但不等价。
\item
  设  为向 2-进域 \(K_{0}\) 添加所有多项式 \(Q_{2}\) 的根所得的域；令 v 为 K＿0 上延拓 2-进赋值的唯一赋值，并令 K 为 K＿0 关于 v 的完备化。证明多项式 \(X^{2^{n}} - 2\) 在 K\(K_{0}\)\(K_{0}\)\(X^{2} - 3\){[}X{]} 中不可约；令  为该多项式的根域，令 \(K'\) 为 v 到 \(v'\) 的延拓；则\(K'\) 的延拓；则
\end{enumerate}

\[
n = \left[ \mathrm{K} ^ {\prime}: \mathrm{K} \right] = 2, \quad e (v ^ {\prime} / v) = f (v ^ {\prime} / v) = 1.
\]

\begin{enumerate}
\def\labelenumi{\arabic{enumi}.}
\setcounter{enumi}{2}
\tightlist
\item
  令 k 为素域  上无穷多个不定元的有理函数域 ，并令 \(\mathbf{F}_{p}(\mathbf{X}_{n})_{n\in\mathbf{N}}\) 为 k 上两个不定元的有理函数域。\(F_{p}\)\(K = k(\mathbf{U}, \mathbf{V})\) 为 k 上两个不定元的有理函数域。
\end{enumerate}

\begin{enumerate}
\def\labelenumi{(\alph{enumi})}
\item
  证明形式幂级数域  的元素  在有理函数域 \(\mathbf{P}(\mathbf{U}) = \sum_{n=0}^{\infty} \mathbf{X}_n^p \mathbf{U}^{np}\) 上不代数。映射 \(k((\mathbf{U}))\)\(k(\mathbf{U})\) 从 \(\mathbf{F}(\mathbf{U}, \mathbf{V}) \mapsto \mathbf{F}(\mathbf{U}, \mathbf{P}(\mathbf{U}))\) 到 \(k[\mathbf{U}, \mathbf{V}]\) 可延拓为 \(k((\mathbf{U}))\) 到 \(\mathbf{K}\) 的一个子域上的同构；将 \(k((\mathbf{U}))\) 上等于形式幂级数的阶的赋值（§ 3，第 3 节，例 3）限制到该子域，得到 \(k((\mathbf{U}))\) 上的离散赋值 v，其剩余域为 k。\(v\)\(\mathbf{K}\)\(k\) 上的离散赋值 v，其剩余域为 k。
\item
  令  为  的代数扩张，使得 \(\mathbf{K}'\)\(\mathbf{K}(\mathbf{V}^{1/p})\)\(\mathbf{K}\)；若 \([\mathbf{K}':\mathbf{K}] = p\) 为 \(v'\) 上延拓 v 的唯一赋值，证明 \(\mathbf{K}'\)。因此，赋值 \(v\)\(e(v'/v) = f(v'/v) = 1\) 的环不是赋值 v 的环上的有限生成模。\(v'\)\(v\) 的环不是赋值 v 的环上的有限生成模。
\end{enumerate}

\begin{enumerate}
\def\labelenumi{\arabic{enumi}.}
\setcounter{enumi}{3}
\item
  设  为域，\(k\) 为 k 上两个不定元的有理函数域，v 为 \(\mathbf{K} = k(\mathbf{X}, \mathbf{Y})\) 上取值于按字典序排序的群 \(k\)\(v\)\(\mathbf{K}\) 的赋值，满足 \(\mathbf{Z} \times \mathbf{Z}\)、\(v(\mathbf{X}) = (0, 1)\)。令 \(v(\mathbf{Y}) = (1, 0)\) 为域 \(\mathbf{K}'\)；证明 v 到 \(\mathbf{K}(\sqrt{\mathbf{X}})\) 有唯一延拓 ，且 \(v\)\(v'\)、\(\mathbf{K}'\)\(e(v'/v) = 2\)，但赋值 \(f(v^{\prime}/v)=1\) 的环不是赋值 v 的环上的有限生成模。\(v^{\prime}\)\(v\) 的环不是赋值 v 的环上的有限生成模。
\item
  设 K 为域，A 为 K 的赋值环，L 为 K 的有限代数扩张。令  为 K 的另一个赋值环；令 \(A' \supset A\)（）为 L 中满足  的赋值环，并令 \(A_i' (1 \leqslant i \leqslant m)\)\(A_i' \cap K = A'\) 为它们各自的剩余域。令 k 为 \(k_i'\) 的剩余域，\(A'\) 为 A 的典范像所在的 K 的赋值环；最后，令 \(A''\)（）为  的满足 \(A_{ij}' (1 \leqslant j \leqslant n_i)\)\(k_i'\) 的赋值环，并令 \(A_{ij}'' \cap k = A''\) 为 \(A_{ij}\) 中 \(A_i'\) 的逆像（）。\(A_{ij}'' (1 \leqslant i \leqslant m, 1 \leqslant j \leqslant n_i)\)）。
\end{enumerate}

\begin{enumerate}
\def\labelenumi{(\alph{enumi})}
\item
  证明  是 \(\mathbf{A}_{ij}\) 的彼此不同的赋值环，满足 \(\mathbf{L}\)，并且 \(\mathbf{A}_{ij} \cap \mathbf{K} = \mathbf{A}\) 的每个满足 \(\mathbf{B}\)\(\mathbf{L}\) 的赋值环 \(\mathbf{B} \cap \mathbf{K} = \mathbf{A}\) 都等于某个 。\(\mathbf{A}_{ij}\)。
\item
  证明
\end{enumerate}

\[
e \left(\mathrm{A} _ {i j} / \mathrm{A}\right) = e \left(\mathrm{A} _ {i} ^ {\prime} / \mathrm{A} ^ {\prime}\right) e \left(\mathrm{A} _ {i j} ^ {\prime \prime} / \mathrm{A} ^ {\prime \prime}\right) \quad \text { and } \quad f \left(\mathrm{A} _ {i j} / \mathrm{A}\right) = f \left(\mathrm{A} _ {i j} ^ {\prime \prime} / \mathrm{A} ^ {\prime \prime}\right).
\]

¶ 6. 设 K 为域，v 为 K 上的赋值，A 为 v 的环。

\begin{enumerate}
\def\labelenumi{(\alph{enumi})}
\item
  假设 A 是 Hensel 的（第三章 § 4，习题 3）；此时我们也约定称 K 关于 v 是 Hensel 的。于是对 K 的每个代数扩张 L 以及 L 上延拓 v 的每个赋值 ，\(v'\) 的环  都是 Hensel 的（第三章 § 4，习题 4）。\(A'\)\(v'\) 都是 Hensel 的（第三章 § 4，习题 4）。
\item
  若 K 关于 v 完备且 v 的高度为 1，则 A 是 Hensel 的。给出 K 关于 v 完备、v 的高度为 2 而 A 不是 Hensel 环的例子（参见习题 1）。
\item
  若 K 关于 v 线性紧，证明 K 是 Hensel 的。（在第三章 § 4，习题 3 的条件 (H) 的记号中，令 L 表示 A 的素理想所成的集合（按包含关系）全序，令 \(\mathcal{L}\) 表示有序对 (B, φ) 所成的集合，其中 B 是 L 的良序子集（关于关系 ⊃），而 φ: p ↦ (Q\_p, Q'\_p) 是从 B 到集合 A{[}X{]} × A{[}X{]} 的映射，满足以下性质：(1) Q\_p 与 Q'\_p 分别为次数 deg( \(\overline{Q}\) ) 与 deg( \(\overline{Q}'\) ) 的首一多项式；(2) 若 p ⊃ q 是 B 的两个元素，则多项式 Q\_p - Q\_q 与 Q'\_p - Q'\_q 的系数属于 p；(3) 多项式 P - Q\_pQ'\_p 的系数属于 p；(4) \(\bar{f}(Q_{p}) = \overline{Q}, \bar{f}(Q'_{p}) = \overline{Q}'\)。在 L 上规定序：当 B ⊂ B' 且 φ' 是 φ 到 B' 的延拓时，置 (B, φ) ≤ (B', φ')。证明 L 有极大元 (B₀, φ₀)，且 B₀ 有一个等于 (0) 的末元：反设之，并根据 B₀ 是否有末元分两种情形。前一种情形中，若 p 是该末元，取 c ∈ A，使得 v(c) 是 v 在多项式 P - Q\_pQ'\_p 的系数集上所取的最小值，并令 p' 为包含 c 的 A 的最小素理想；若 p' = p，则使用 § 4 的习题 2 (a) 以及线性紧模被闭子模取商后仍线性紧、因而完备这一事实。反之，若 B₀ 没有末元，则直接使用 A 线性紧的假设。） 
\item
  用习题 5 的记号，证明 A 为 Hensel 环的充要条件是  与 \(\mathbf{A}'\) 都是 Hensel 环。（注意，若 \(\mathbf{A}''\)，则存在 \(\mathbf{P} \in \mathbf{A}'[\mathbf{X}]\)，使得 \(a \in \mathbf{A}\)（其中 n 是 \(a^n\mathbf{P}(\mathbf{X}/a) \in \mathbf{A}[\mathbf{X}]\) 的次数）；为验证 \(n\)\(\mathbf{P}\) 的第三章 § 4，习题 3 的公理 (H)，只需考虑 \(\mathbf{A}'\) 中的多项式；使用 \(\mathbf{A}[\mathbf{X}]\) 这一事实，其中 \(\mathbf{A}' = \mathbf{A}_{\mathfrak{p}}\) 是包含于极大理想 \(\mathfrak{p}\) 的  的素理想，并注意若 \(\mathbf{A}\) 中的两个多项式强互素，则它们在 \(\mathfrak{m}\)\((\mathbf{A}/\mathfrak{p})[\mathbf{X}]\) 中的像也强互素。）\((\mathbf{A}/\mathfrak{m})[\mathbf{X}]\) 中的像也强互素。）
\item
  假设 A 是 Hensel 环；于是对 K 的每个代数扩张 L，在 L 上存在（除等价外）唯一一个延拓 v 的赋值。（若  是首一不可约多项式，\(P \in A[X]\)（\(x_{i}\)）是 P 在其分裂域 N 中的不同根，证明对 N 上每个延拓 v 的赋值 \(1 \leqslant i \leqslant m\)，当 \(v'\) 时各  都相等。）反之，若赋值环 A 具有此性质，则 A 是 Hensel 环（注意在 A 在 L 中的整闭包 \(v'(x_{i})\) 中，位于 m 之上的极大理想至多有一个（第五章 § 1，习题 13），从而若 \(1 \leqslant i \leqslant m\)\(A'\) 是首一不可约多项式，则其在 \(P \in A[X]\) 中的像不可能是两个互素多项式的乘积）。\((\mathbf{A}/\mathfrak{m})[\mathbf{X}]\) 中的像不可能是两个互素多项式的乘积）。
\end{enumerate}

\begin{enumerate}
\def\labelenumi{\arabic{enumi}.}
\setcounter{enumi}{6}
\tightlist
\item
  设 K 为域，v 为 K 上的赋值，A 为 v 的环，m 为 v 的理想。假设 A 是 Hensel 环；令 L 为 K 的 n 次有限扩张；令 \(v'\) 为 L 上延拓 v 的唯一赋值（习题 6 (e)），A' 为其环，\(m'\) 为其理想。令 x 为 A' 的任意元素；证明 x 在 A'/m' 中的类 \(\bar{x}\) 在 A/m 上的次数整除 n，且 \(v'(x)\) 在 \(\Gamma_{v'}/\Gamma_v\) 中的类的阶也整除 n（考虑 x 在 K 上的极小多项式）。
\end{enumerate}

¶ 8. 设 K' 为域，B 为 K' 上赋值 v 的环，G 为 K' 的有限自同构群，K 为 K' 中由 G 不变元素组成的子域，且 A = K ∩ B 是 K 的赋值环，对应于赋值 w = v\textbar K。

\begin{enumerate}
\def\labelenumi{(\alph{enumi})}
\item
   上延拓 \(\mathbf{K}'\) 的赋值是 \(w\)（其中 \(v \circ \sigma\)）（第 6 节，命题 6 的推论 1），而 \(\sigma \in \mathcal{G}\) 在 \(\mathbf{A}'\)\(\mathbf{A}\) 中的整闭包 \(\mathbf{K}'\) 是所有  的交，其中 \(\sigma \cdot \mathbf{B}\) 遍历 \(\sigma\)。若 \(\mathcal{G}\) 是 \(\mathfrak{p}(\mathbf{B})\) 与极大理想 \(\mathbf{A}'\) 的交，则 \(\mathfrak{m}(\mathbf{B})\)。证明分解群 \(\sigma \cdot \mathfrak{p}(\mathbf{B}) = \mathfrak{p}(\sigma \cdot \mathbf{B})\)（第五章 § 2，第 2 节）是由满足 \(\mathcal{G}^{\mathbb{Z}}(\mathfrak{p}(\mathbf{B}))\) 的  组成的  的子群；记 \(\mathcal{G}\)，并令 \(\sigma\)\(\sigma \cdot \mathbf{B} = \mathbf{B}\)\(\mathcal{G}^{\mathbb{Z}} = \mathcal{G}^{\mathbb{Z}}(\mathfrak{p}(\mathbf{B}))\) 为 \(\mathbf{K}^{\mathbb{Z}}\) 中由 \(\mathbf{K}'\) 不变元素组成的子域，令 \(\mathcal{G}^{\mathbb{Z}}\) 为 v 在 \(\mathbf{B}^{\mathbb{Z}}\) 上诱导的赋值的环，它等于 \(\mathbf{K}^{\mathbb{Z}}\)；回忆 \(v\)\(\mathbf{B} \cap \mathbf{K}^{\mathbb{Z}}\) 与 \(\mathbf{B}^{\mathbb{Z}}\) 的剩余域相同，且极大理想 \(\mathbf{A}\) 由 \(\mathfrak{m}(\mathbf{B}^{\mathbb{Z}})\) 生成（第五章 § 2，第 2 节，命题 4）。\(\mathfrak{m}(\mathbf{A})\mathbf{B}^{\mathbb{Z}}\) 生成（第五章 § 2，第 2 节，命题 4）。
\item
  令  表示 v 到 \(v^{\mathbf{Z}}\) 的限制，令 \(v\)\(\mathbf{K}^{\mathbf{Z}}\) 与 \(\Gamma\) 表示 w 与 \(\Gamma^{\mathbf{Z}}\) 的阶群。证明 \(w\)\(v^{\mathbf{Z}}\)。（令 \(\Gamma = \Gamma^{\mathbf{Z}}\) 为 \(\mathcal{G}^{\mathrm{D}} \supset \mathcal{G}^{\mathbf{Z}}\) 的子群，它由满足 \(\mathcal{G}\) . B 依赖于 B（§ 7，第 2 节）的  组成。对 \(\sigma \in \mathcal{G}\) 的阶数作归纳。若 \(\sigma\)\(\mathcal{G}\)，令 \((\mathcal{G}: \mathcal{G}^{\mathrm{D}}) > 1\) 为 \(\mathbf{K}''\) 的不变域，令 \(\mathcal{G}^{\mathrm{D}}\)；证明 \(\mathbf{A}'' = \mathbf{K}'' \cap \mathbf{B}\) 的阶群等于 \(v|\mathbf{K}''\)：为此使用逼近定理（§ 7，第 2 节，定理 1），证明对所有 \(\Gamma\)，存在 \(x \in K''\)，使得 \(y \in K''\)，且相对于 K 的 y 的共轭元中除 y 外的各个 \(v(y) = v(x)\) 满足 ；于是可将 G 换为 \(v(y_i) = 0\)，并应用归纳假设。若 \(y_i\)\(y\)\(K\)\(y\)\(G\)\(G^D\)，令 \(G^D = G\) 为由 \(B_0\)（其中 \(K'\)\(\sigma.B\)）生成的 \(\sigma \in G\) 的赋值环；令  为其剩余域，\(\bar{K}' = \kappa(B_0)\) 为典范同态，\(\pi: B_0 \to \bar{K}'\) 为 \(\bar{B} = \pi(B)\) 的赋值环，\(\bar{K}'\) 为相应的 \(\bar{v}\) 上的赋值，使得在 \(\bar{K}'\) 中 ，从而 \(v = \bar{v} \circ \pi\)（\(B_0\)\(\bar{\Delta}\) 的阶群）是 \(\bar{v}\)（v 的阶群）的子群；若 \(\Delta\)，\(v\)\(\Delta_0 = \Delta/\bar{\Delta}\) 为典范同态，则 \(\bar{\omega}: \Delta \to \Delta_0\) 是对应于 \(v_0 = \bar{\omega} \circ v\) 且在 G 下不变的  上的赋值。取商后，G 在 \(K'\) 上定义一个自同构群 ；令 N 为典范同态 \(B_0\)\(G\)\(G\)\(\bar{G}\) 的核；证明 \(\bar{K}'\)\(N\)\(G \to \bar{G}\)，并推出若 \(N \subset G^Z\)，可将 G 换为 \(N \neq \{e\}\)，然后应用归纳假设。最后假设 \(G\)\(G/N\)；令 \(N = \{e\}\)，令 \(A_0 = K \cap B_0\) 为 \(w_0\) 到 K 的限制；证明 \(v_0\) 的阶群为 \(K\)\(w_0\)，且 \(\Delta_0\) 是 \(\pi(A_0) = \bar{K}\) 的不变域（参见第 1 节，引理 2）；若 \(\bar{G}\) 为 \(\bar{w}\) 到  的限制，\(\bar{K}\) 为其阶群，则 \(\bar{v}\)\(\bar{\Gamma}\) 典范同构于 \(\Delta_0\)。另一方面，群 \(\Gamma/\bar{\Gamma}\) 是 \(G^Z(p(\bar{B}))\) 在 \(G^Z(p(B))\) 中的典范像；注意 \(G^Z\)，因而可应用归纳假设证明 \(G^D \neq G^Z\) 等于 \(\bar{\Gamma}\) 到 \(F^Z(θ)\)\(\bar{v}\) 的限制的阶群 \(K^Z\) 的限制的阶群 
\item
  令  为  的惯性群（第五章 § 2，第 2 节），并令  表示 \(\mathcal{G}^{\mathrm{T}} = \mathcal{G}^{\mathrm{T}}(\mathfrak{p}(\mathrm{B}))\)\(\mathfrak{p}(\mathrm{B})\)\(\mathbf{K}^{\mathrm{T}}\) 中由 \(\mathbf{K}'\) 的不变元素组成的子域，令 \(\mathcal{G}^{\mathrm{T}}\) 表示 v 在 \(\mathbf{B}^{\mathrm{T}}\) 上诱导的赋值的环，它等于 \(\mathbf{K}^{\mathrm{T}}\)；回忆：若 \(v\)\(\mathbf{B} \cap \mathbf{K}^{\mathrm{T}}\) 与 \(k\) 分别为 A 与 B 的剩余域，则 \(k'\) 的剩余域  是 k 的拟 Galois 扩张 \(k^{\mathrm{T}}\) 中的最大可分扩张，并且 \(\mathbf{B}^{\mathrm{T}}\)\(k\)\(k'\) 典范同构于 \(k\)\(\mathcal{G}^{\mathrm{Z}} / \mathcal{G}^{\mathrm{T}}\)（或 \(k\)\(k'\)）的 \(k^{\mathrm{T}}\) -自同构群（第五章 § 2，定理 2 与命题 5）。令  为 v 到 \(v^{\mathrm{T}}\) 的限制；证明 \(v\)\(\mathbf{K}^{\mathrm{T}}\) 的阶群也等于 \(v^{\mathrm{T}}\)（将引理 2（第 1 节）应用于 \(\Gamma\) 对 \(\mathbf{K}^{\mathrm{T}}\) 的扩张）。\(\mathbf{K}^{\mathrm{Z}}\) 的扩张）。
\end{enumerate}

\begin{enumerate}
\def\labelenumi{\arabic{enumi}.}
\setcounter{enumi}{8}
\tightlist
\item
  设 K 为域，L 为 K 的有限代数扩张，\(v'\) 为 L 上的赋值，A' 为其环，v 为 \(v'\) 到 K 的限制，\(A = A' \cap K\) 为其环。
\end{enumerate}

\begin{enumerate}
\def\labelenumi{(\alph{enumi})}
\item
  假设 A 是 Hensel 环（习题 6）；证明乘积  整除 \(e(v'/v)f(v'/v)\)，且商 \(n = [\mathbf{L}:\mathbf{K}]\) 是 v 的剩余域 k 的特征指数的幂。（将其归结为 L 是 K 的 Galois 扩张的情形；使用习题 6 (e)、习题 7 以及第五章 §2 的定理 2 与命题 5。）\(n / e(v'/v)f(v'/v)\)\(k\)\(v\) 是 v 的剩余域 k 的特征指数的幂。（将其归结为 L 是 K 的 Galois 扩张的情形；使用习题 6 (e)、习题 7 以及第五章 §2 的定理 2 与命题 5。）
\item
  进一步假设 n 不被 v 的剩余域 k 的特征指数整除，且 。证明：对于群 \(n\)\(k\)\(v\)\(f(v' / v) = 1\) 中元素  的阶所等于的每个整数 m，都存在 \(m\)\(\gamma\)，使得 \(\Gamma_{v'} / \Gamma_v\)\(x \in \mathbf{L}\) 的类等于 \(v'(x) \mod \Gamma_v\)，且 \(\gamma\)。\(x^m \in \mathbf{K}\)。
\end{enumerate}

¶ 10. 令 w 为 2-进域  上的 2-进赋值；在有理函数域 \(Q_{2}\) 上考虑延拓 w 且满足\(\mathbf{Q}_2(\mathbf{X})\)\(v\)\(w\) 上考虑延拓 w 且满足

\[
v (a _ {0} + a _ {1} \mathrm{X} + \dots + a _ {n} \mathrm{X} ^ {n}) = \inf _ {i} (w (a _ {i}))
\]

（§ 10，第 1 节，引理 1）；令 K 为  关于此赋值的完备化，并仍以 v 表示其赋值。令 L 为多项式 \(\mathbf{Q}_2(\mathbf{X})\) 在 K\(v\)\((\mathrm{Y}^2 - \mathrm{X})^2 - 2\){[}Y{]} 中的分裂域，令 \(v'\)\(v\) 为 L 上延拓 v 的唯一赋值。证明 {[}L: K{]} = 8，，\(e(v'/v) = 4\)；若 \(f(v'/v) = 2\)（分别 \(k\)）是 v（分别 \(k'\)）的剩余域，则 \(v\)\(v'\) 是 k 的二次根式扩张；证明不存在 L 的子扩张 E 满足 \(k'\)\(k\){[}E: K{]} = 2 且  到 E 的限制的剩余域等于 \(v'\)。（必然有 \(k'\)，其中 \(\mathrm{E} = \mathrm{K}(\sqrt{\alpha})\) 不是 K 中的平方，但满足 \(\alpha \in \mathrm{K}\)（模 2）；借助 L 在 \(\alpha \equiv \mathrm{X}\) 上的适当基表示 ，并注意 \(\sqrt{\alpha}\) 中不存在平方同余于 X 模 2 的元素。）\(\mathrm{K}(\sqrt{2})\)\(\mathrm{K}(\sqrt{2})\) 中不存在平方同余于 X 模 2 的元素。）

¶ 11. 设 K 为域，L 为 K 的有限 Galois 扩张，G 为其 Galois 群。令 v 为 L 上的赋值，其剩余域为 k、阶群为 Γ；假设 v 在 G 下不变，且 v\textbar K 的限制与 v 具有相同的剩余域 k，于是 \(G = G^{Z} = G^{T}\)（习题 8 的记号）。

\begin{enumerate}
\def\labelenumi{(\alph{enumi})}
\item
  令 、\(x \in \mathbf{L}^*\)。证明 \(\sigma \in \mathcal{G}\) 是 v 的单位，且它在 \(x^{-1}\sigma(x)\) 中的像  只依赖于赋值 \(v\)\(\varepsilon_{\sigma}(x)\) 以及 \(k^*\)\(v(x)\) 模 \(\sigma\) 的换位子群  的类，因而定义了一个 Z-双线性映射（也记作 \(\mathcal{G}'\)）\(\mathcal{G}\)\(\varepsilon\)，它在 \((\mathcal{G}/\mathcal{G}') \times \Gamma \to k^*\) 上等于 1。\((\mathcal{G}/\mathcal{G}') \times v(\mathbf{K}^*)\) 上等于 1。
\item
  对所有 ，令 \(\sigma \in \mathcal{G}\) 为从 \(\phi(\sigma)\) 到  的同态，它将  映到 \(\Gamma / v(\mathbf{K}^*) \to k^*\)\(\varepsilon_{\sigma}(x)\)（对所有 \(v(x)\)）；于是定义了同态\(x \in \mathbf{L}^*\)）；于是定义了同态
\end{enumerate}

\[
\phi \colon \mathcal {G} \to \operatorname{Hom} _ {\mathbf {Z}} (\Gamma / v (\mathrm{K} ^ {*}), k ^ {*}).
\]

证明：若 p 是 k 的特征指数，则  的核  的阶是 p 的幂。（否则， 中存在 \(p\)\(k\)\(\mathcal{N}\) 以及素数 \(\phi\)\(p\)\(\sigma \neq e\)，使得 \(\mathcal{N}\)\(q \neq p\)；对 L - K 中某个整元 x，写成 \(\sigma^q = e\)，其中 \(x\)\(L - K\)\(\sigma(x) = x + y\)，并计算 \(v(y) > v(x)\) 得到矛盾。）推出 \(\sigma^q(x)\) 是可解群。\(\mathcal{G}\) 是可解群。

\begin{enumerate}
\def\labelenumi{(\alph{enumi})}
\setcounter{enumi}{2}
\tightlist
\item
  假设  与 p 互素；证明 \(n = [\mathbf{L} : \mathbf{K}]\) 是双射，\(p\)\(\phi\)，并且若 \(e(\mathbf{L} / \mathbf{K}) = n\) 是 \(\nu\) 中元素的阶的最小公倍数，则 \(lcm\)\(\mathcal{G}\) 含有 \(k^*\) 次单位根。（使用 (b) 及第 1 节的引理 2。）\(\nu\) 次单位根。（使用 (b) 及第 1 节的引理 2。）
\end{enumerate}

¶ 12. 设 K 为域，v 为 K 上的赋值，A 为赋值 v 的环，Γ 为其阶群。

\begin{enumerate}
\def\labelenumi{(\alph{enumi})}
\tightlist
\item
  令 \(f(\mathbf{X}) = a_{0}\mathbf{X}^{n} + a_{1}\mathbf{X}^{n-1} + \cdots + a_{n}\) 为 A{[}X{]} 中满足 \(v(a_{0}) = 0\) 且所有根都属于 K 的多项式；令 \(x_{i} (1 \leqslant i \leqslant r)\) 为这些不同的根，令 \(k_{i}\) 为 \(x_{i} (1 \leqslant i \leqslant r)\) 的重数。令
\end{enumerate}

\[
g (\mathbf {X}) = b _ {0} \mathbf {X} ^ {n} + b _ {1} \mathbf {X} ^ {n - 1} + \dots + b _ {n} = b _ {0} (\mathbf {X} - y _ {1}) \dots (\mathbf {X} - y _ {n})
\]

另一个满足  且  属于 K 的 A{[}X{]} 中的多项式。证明存在 \(v(b_{0}) = 0\)\(y_{h}\)\(\lambda_{0} \in \Gamma\)，使得对每个不同指标 i、j 的有序对都有 \(v(x_{i} - x_{j}) < \lambda_{0}\)，并且对所有 \(\lambda \geqslant \lambda_{0}\)，存在 \(\mu \geqslant \lambda\)，使得对所有 i 满足 \(v(a_{i} - b_{i}) \geqslant \mu\) 时，对 \(1 \leqslant i \leqslant r\)，恰有 \(k_{i}\) 个指标 h 满足 \(v(x_{i} - y_{i}) \geqslant \lambda\)。（先用两种方法计算 \(v(f(y_{h}))\)，证明必有某个 i 使 \(v(x_{i} - y_{h}) \geqslant \lambda\)；然后计算

\[
v \left(\frac {f ^ {\prime} (z)}{f (z)} - \frac {g ^ {\prime} (z)}{g (z)}\right)
\]

在适当的点  处的值。）若进一步除 \(z \in \mathbf{K}\) 外均有 ，证明只要 \(b_{i} = a_{i}\) 充分大，\(i = n - 1\)\(\lambda_0\) 就两两不同（假设 \(y_{h}\) 不是 g 的单根，计算 ）。\(f'(y_h)/f(y_h)\)\(y_{h}\)\(g\)）。

\begin{enumerate}
\def\labelenumi{(\alph{enumi})}
\setcounter{enumi}{1}
\tightlist
\item
  以下假设 \(\mathbf{K}\) 是子域 \(\mathbf{K}_0\) 的代数闭包，且 \(\mathbf{A}_0 = \mathbf{A} \cap \mathbf{K}_0\) 是 Hensel 环。假设多项式 \(f\) 属于 \(\mathbf{A}_0[\mathbf{X}]\)，且在 \(\mathbf{K}_0\) 上不可约且可分。令 \(z \in \mathbf{K}\) 满足
\end{enumerate}

\[
v (z - v _ {i}) > v (z - x _ {j})
\]

对所有 \(j \neq i\) 成立；证明 \(\mathrm{K}_0(x_i) \subset \mathrm{K}_0(z)\)。（证明 \(z - x_i\) 关于 \(\mathrm{K}_0(x_i)\) 的所有共轭元都相等。）

\begin{enumerate}
\def\labelenumi{(\alph{enumi})}
\setcounter{enumi}{2}
\tightlist
\item
  假设  可分，并令 \(f \in \mathbf{A}_0[\mathbf{X}]\) 为 f 在 \(f = a_0 \prod_{i=1}^{\infty} f_i\) 中分解为首一不可约多项式的分解（事实上它们属于 \(f\)\(\mathbf{K}_0[\mathbf{X}]\)，参见第五章 § 1，第 3 节，命题 11）。证明（沿用 (a) 的记号）：若 \(\mathbf{A}_0[\mathbf{X}]\) 取充分大，则 g 分解为首一\(\mu\)\(g\) 取充分大，则 g 分解为首一
\end{enumerate}

可约因子在 \(K_{0}[X]\) 中可写成 \(b_{0}\prod_{i=1}^{n}g_{i}\)，其中对所有 i，\(g_{i}\) 与 \(f_{i}\) 次数相同（称为与 f 同类型的分解），并且对所有 i，\(f_{i}\) 与 \(g_{i}\) 的分裂域相同。（先证明 g 可分；再考虑一个包含 f 与 g 的根的有限 Galois 扩张 \(K_{0}\)，考察该扩张的 Galois 群如何置换这些根；最后使用 (b)。）

\begin{enumerate}
\def\labelenumi{(\alph{enumi})}
\setcounter{enumi}{3}
\tightlist
\item
  给出一个 f 不可约但不可分的例子，并且对所有 ，存在多项式\(f\)\(\mu \geqslant \lambda_0\)，存在多项式
\end{enumerate}

\[
g (\mathbf {X}) = b _ {0} \mathbf {X} ^ {n} + b _ {1} \mathbf {X} ^ {n - 1} + \dots + b _ {n} \in \mathrm{A} _ {0} [ \mathbf {X} ]
\]

使得对所有 i 有 \(v(a_{i} - b_{i}) \geqslant \mu\)，但它不可约（取 \(K_{0}\) 使得 \(K \neq K_{0}\)，且 K 是 \(K_{0}\) 的根式扩张）。

¶ 13. (a) 设 K 为域，v 为 K 上的非真赋值。令 E 为由超滤子 U 滤过的集合，\(\xi \mapsto x_{\xi}\)、\(\xi \mapsto y_{\xi}\) 为 E 到带拓扑 \(\mathcal{T}_{v}\) 的 K 的两个映射。证明：若映射 \(\xi \mapsto x_{\xi}y_{\xi}\) 关于 U 在 K 中收敛于 0，则两个映射 \(\xi \mapsto x_{\xi}\)、\(\xi \mapsto y_{\xi}\) 中至少有一个也如此。（注意，若 \(\xi \mapsto x_{\xi}\) 关于 \(\mathfrak{U}\) 没有聚点，则映射 \(z \mapsto x_{\xi}^{-1}\) 在 \(\mathfrak{U}\) 的某个子集上有界。）

\begin{enumerate}
\def\labelenumi{(\alph{enumi})}
\setcounter{enumi}{1}
\item
  令  为 \(f\) 中的非常数多项式。证明：若 \(\mathbf{K}[\mathbf{X}]\) 关于 \(\xi \mapsto f(x_{\xi})\) 在  中收敛于 0，则 \(\mathbf{K}\) 在 \(\mathfrak{U}\)\(\xi \mapsto x_{\xi}\) 中收敛（将 f 分解为 \(\hat{\mathbf{K}}\) 的代数扩张中的因子，并使用 (a)）。\(f\)\(\mathbf{K}\) 的代数扩张中的因子，并使用 (a)）。
\item
  假设 K 在  中代数闭。则 K 到自身的映射 \(\hat{\mathbf{K}}\) 是闭映射（使用 (b)）。推出：若 \(x \mapsto f(x)\) 是 \(g\) 中的有理函数，B 是 K 的有界闭子集，则 \(\mathbf{K}(\mathbf{X})\)（其中 P 是 g 在 B 中的极点集）在 K 中闭。\(g(\mathbf{B} - \mathbf{P})\)\(g\)（其中 P 是 g 在 B 中的极点集）在 K 中闭。
\item
  假设 \(\mathbf{K}\) 的代数闭包是 \(\mathbf{K}\) 的根式扩张。证明 \(\hat{\mathbf{K}}\) 是代数闭域。（将 (b) 应用于 \(\hat{\mathbf{K}}\) 的代数闭包及多项式 \(f \in \hat{\mathbf{K}}[\mathbf{X}]\)；还要使用习题 12 (a)。）
\end{enumerate}

¶ 14. (a) 域 K 关于赋值 v 为 Hensel 的充要条件是：K 是包含于  中的 K 的最大可分代数扩张，且 \(\hat{K}\) 是 Hensel 的。（为看出条件必要，使用习题 12 (b)；为证明充分，注意为验证第三章 § 4，习题 3 的条件 (H)，只需考虑 P 在 K 上可分的情形；因为若 Q 是 P 在 \(\hat{K}\) 中的不可约因子，且 \(\hat{K}[X]\) 属于 K\(Q^{p^{e}}\){[}X{]}，则 Q 是 \(Q^{p^{e}}\) 中 P 与  的最大公因子。）考虑 v 的高度为 1 的情形（参见习题 6 (b)）。\(\hat{K}[X]\) 的最大公因子。）考虑 v 的高度为 1 的情形（参见习题 6 (b)）。

\begin{enumerate}
\def\labelenumi{(\alph{enumi})}
\setcounter{enumi}{1}
\item
  由 (a) 推出 p-进域  中存在一个关于 p-进赋值为 Hensel 的可数子域。\(p\)\(\mathbf{Q}_p\)\(p\) 中存在一个关于 p-进赋值为 Hensel 的可数子域。
\item
  给出一个关于离散赋值为 Hensel 但不完备的域 K，使得 \(\hat{K}\) 是 K 的有限根式扩张（参见第五章 § 1，习题 20）。
\end{enumerate}

¶ 15. (a) 设 K 为域，\(v_1\)、\(v_2\) 为 K 上的两个独立赋值（§ 7，第 2 节），\(K_1\)、\(K_2\) 分别为 K 关于 \(v_1\)、\(v_2\) 的完备化。令 \(L_1\)、\(L_2\) 为满足 \(K \subset L_1 \subset K_1\)、\(K \subset L_2 \subset K_2\) 的两个域，它们是 Hensel 的（分别通过 \(v_1\)、\(v_2\) 的连续延拓）。令 \(g_1 \in L_1[X]\)、\(g_2 \in L_2[X]\) 为次数相同的两个可分多项式，次数为 n。证明存在多项式 \(h \in K[X]\)，使得 h 在 \(L_i[X]\) 中的分解类型（习题 12 (c)）与 \(g_i\) 相同（\(i = 1, 2\)）。（将其归结为 \(g_1\)、\(g_2\) 属于 K{[}X{]} 且为首一的情形；考虑多项式

\[
a ^ {n} g _ {1} (\mathbf {X} / a) + b ^ {n} g _ {2} (\mathbf {X} / a) - \mathbf {X} ^ {n},
\]

其中 \(v_{1}(a) \leqslant 0\)，\(v_{2}(a) > 0\) 任意大，且 \(v_{1}(b) > 0\) 任意大。）

\begin{enumerate}
\def\labelenumi{(\alph{enumi})}
\setcounter{enumi}{1}
\item
  由 (a) 推出：若 K 关于 \(v_{1}\) 为 Hensel，则 \(L_{2}\) 的代数闭包是 \(L_{2}\) 的根式扩张，且 \(K_{2}\) 是代数闭域（令 \(g_{2}\) 不可约，令 \(g_{1} \in K[X]\) 为 n 个不同的素次数因子的乘积）。
\item
  由 (b) 推出：若 \(\mathbf{K}\) 关于 \(v_{1}\) 与 \(v_{2}\) 都为 Hensel，则 \(\mathbf{K}\) 的代数闭包是 \(\mathbf{K}\) 的根式扩张。
\item
  证明若 K 关于离散赋值 v 为 Hensel，则 K 的代数闭包不可能是 K 的根式扩张，因而 K 不可能关于任何独立于 v 的赋值为 Hensel（使用代数，第五章 § 11，习题 12）。
\end{enumerate}

\begin{enumerate}
\def\labelenumi{\arabic{enumi}.}
\setcounter{enumi}{15}
\tightlist
\item
  设 K 为带有高度 1 赋值 v 的 Hensel 域。若 L 是 K 的无限可分代数扩张，证明 L 不可能关于延拓 v 的赋值（也记作 v）完备。（构造 L 的元素序列 ，使得 \((x_{p})\) 关于 K 的次数  趋于 \(n_{p}\)，且 \(x_{p}\)\(+\infty\) 严格大于 \(v(x_{p+1}-x_{p})\) 与其关于 K 的共轭元之间的差的赋值；证明序列 \(x_{p}\) 是不能在 E 中收敛的 Cauchy 序列，使用习题 12 (b)。）\((x_{p})\) 是不能在 E 中收敛的 Cauchy 序列，使用习题 12 (b)。）
\end{enumerate}

¶ 17. 设 K 为带有一个非代数闭赋值 v 的 Hensel 域，L 为 K 的满足 {[}K: L{]} 有限的子域。证明在这些条件下 L 关于 v 的限制是 Hensel 的。（反设不然，则存在 v|L 到 K 的代数闭包  的两个延拓 ，使得 \(v_{1}, v_{2}\)、\(v|L\)\(\Omega\)\(v_{1}\) 到 L 的某个有限代数扩张 L' 的限制不等价。推出在包含 K 的 L 的某个有限拟 Galois 扩张 K' 上存在两个不等价的赋值 \(v_{2}\)，且 K' 关于它们是 Hensel 的；使用习题 15 (c)，证明若 \(L'\)\(K'\)\(v_{1}', v_{2}'\)（p 是 \(K'\)\(E = L^{p^{-\infty}}\) 的特征指数），则 \(\Omega\)，但 \(E \neq \Omega\) 是 E 的有限扩张；借助代数，第六章 § 2，习题 31 完成论证。）\(\Omega\) 是 E 的有限扩张；借助代数，第六章 § 2，习题 31 完成论证。）

¶ 18. 设 K 为带有赋值 v 的 Hensel 域，k 为 v 的剩余域，并假设 k 的每个有限代数扩张都是 k 上的循环扩张（例如 k 有限时即如此）。令 A 为 v 的环，令 \(f \in A[X]\) 为首一多项式，并假设当 \(f(\mathbf{X}) = (\mathbf{X} - \alpha_{1}) \ldots (\mathbf{X} - \alpha_{n})\)（其中 \(\alpha_{i}\) 属于 K 的代数闭包）时，A 中的元素 \(D = \prod_{i < j} (\alpha_{i} - \alpha_{j})^{2}\)（f 的“判别式”）满足 \(v(D) = 0\)。令 \(\bar{f}_{j} (1 \leqslant j \leqslant s)\) 为 f 在 k\(\bar{f}\){[}X{]} 中的典范像  的不可约因子，令 \(r_{j}\) 为 \(\bar{f}_{j}\) 的次数。证明 f 的分裂域在 K 上的 Galois 群（看作 \(\alpha_{i}\) 的置换群）由一个置换 \(\sigma\) 生成，该置换分解为长度分别为 \(r_{1}, r_{2}, \ldots, r_{s}\) 的循环（注意 k 的给定次数扩张除 k-同构外唯一）。推出 D 是 A 中平方的充要条件是 n - s 为偶数（“Stickelberger 定理”；考察 D 在 \(\sigma\) 下不变的条件）。

\begin{enumerate}
\def\labelenumi{\arabic{enumi}.}
\setcounter{enumi}{18}
\tightlist
\item
  设 K 为带有赋值 v 的 Hensel 域。令 E 为 K 上的有限维向量空间，Q 为 E 上的非退化二次型，且关系 \(Q(x) = 0\) 蕴含 x = 0。令
\end{enumerate}

\[
\Phi (x, y) = \mathrm{Q} (x + y) - \mathrm{Q} (x) - \mathrm{Q} (y)
\]

为相伴的双线性型。证明 \(2v(\Phi(x,y)) \geqslant v(Q(x)) + v(Q(y))\)；推出

\[
v (\mathrm{Q} (x + y)) \geqslant \inf (v (\mathrm{Q} (x)), v (\mathrm{Q} (y))).
\]

¶ 20. 设 K 为特征 ≠2 的域，关于赋值 v 为 Hensel；令 ξ ↦ ξ̄ 为 K 的满足 v(ξ̄) = v(ξ) 的对合自同构，令 Φ 为 K 上有限维向量空间 E 上的非退化 Hermite 型。

\begin{enumerate}
\def\labelenumi{(\alph{enumi})}
\item
  令  为 \(U = (\alpha_{ij})\) 关于 E 的一组基 \(\Phi\) 的矩阵；证明存在 v 的阶群中的元素 \((e_i)\)，使得若 \(\lambda\) 是 E 上的另一个 Hermite 型，其关于 \(v\)\(\Phi'\) 的矩阵  对每个有序对 \(U' = (\alpha_{ij}')\) 都满足 ，则 \((e_i)\)\(v(\alpha_{ij} - \alpha_{ij}') > \lambda\) 与 \((i,j)\)\(\Phi\) 等价。（仿照代数第九章 § 6，习题 6，并使用上面的习题 14 (a)。）\(\Phi'\) 等价。（仿照代数第九章 § 6，习题 6，并使用上面的习题 14 (a)。）
\item
  假设  是恒等映射（因而 \(\xi \mapsto \overline{\xi}\) 是对称型），且赋值 \(\Phi\) 离散、赋范并满足 \(v\)；令 \(v(2) = 0\) 为 \(\pi\) 的统一化元。证明 \(v\) 除等价外由其指标 \(\Phi\) 以及分别定义在向量空间 \(\nu\)、\(\Psi_1, \Psi_2\)\(k^r\) 上的两个指标为 0 的对称双线性型 \(k^s\) 所刻画，其中  为 \(k\) 的剩余域，且 \(v\)。（借助 Witt 分解将问题归结为 \(r + s = n - 2\) 的指标为 0 的情形；借助习题 19，证明对所有 \(\Phi\)，由满足 \(i \geqslant 0\) 的  组成的集合  是 v 的环 A 上的模。证明若 \(\mathbf{M}_i\)、\(x \in \mathbf{E}\)\(v(\Phi(x, x)) \geqslant i\)\(v\)\(x \in \mathbf{M}_i\)，则 \(y \in \mathbf{M}_{i+1}\)，因而取商后可由 \(v(\Phi(x, y)) \geqslant i + 1\) 在向量 \(\Phi\)-空间 \(k\) 与 \(\mathbf{M}_0 / \mathbf{M}_1\) 上导出对称双线性型。使用 §3 的习题 6 (c) 以及方程 \(\mathbf{M}_1 / \mathbf{M}_2\) 在 \(\xi^2 = \alpha\) 时在 A 中有解这一事实。）考虑 k 有限的情形（参见代数第九章 §6，习题 4）。\(\alpha \equiv 1 (\text{mod. } \pi)\)\(k\) 时在 A 中有解这一事实。）考虑 k 有限的情形（参见代数第九章 §6，习题 4）。
\end{enumerate}

\begin{enumerate}
\def\labelenumi{\arabic{enumi}.}
\setcounter{enumi}{20}
\tightlist
\item
  设 K 为域，v 为 K 上的离散赋值，A 为 v 的环，\(\pi\) 为 v 的统一化元，k 为 v 的剩余域。
\end{enumerate}

\begin{enumerate}
\def\labelenumi{(\alph{enumi})}
\item
  令 P、R 为 A{[}X{]} 中的两个多项式，P 为首一，并假设：(1) ，其中 \(\deg(R) < h \cdot \deg(P)\) 为整数；(2) P 在 k\(h \geqslant 1\)\(\overline{P}\){[}X{]} 中的典范像  不可约。证明若多项式  在 \(Q = P^{h} + \pi R\) 中可约，则必有 h \(\widehat{K}[X]\)\textgreater{} 1，且 R 在 k\(\overline{P}\)\(\overline{R}\){[}X{]} 中的典范像  被  整除。推出：给定首一不可约多项式  及整数 \(r(X) \in k[X]\)，存在首一不可约可分多项式 \(h \geqslant 1\)，使得其典范像 \(Q \in A[X]\) 等于 \(\overline{Q}\)。\(r^{h}\)。
\item
  设  为 k 的次数 m 的代数扩张，h 为整数 \(k' = k(\alpha)\)。证明存在 K 的次数 hm 的代数扩张 L，使得 L 上除等价外只有一个延拓 v 的赋值 \(k\)\(m\)\(h\)\(\geqslant 1\)，且 \(L\)\(K\)\(hm\)\(v'\)、\(L\)\(v\)\(e(v'/v) = h\)，并且 \(f(v'/v) = m\) 的剩余域同构于 \(v'\)（使用 (a)）。\(k'\)（使用 (a)）。
\end{enumerate}

\begin{enumerate}
\def\labelenumi{\arabic{enumi}.}
\setcounter{enumi}{21}
\tightlist
\item
  \begin{enumerate}
  \def\labelenumii{(\alph{enumii})}
  \tightlist
  \item
    若域 \(\mathbf{K}\) 上存在离散赋值，证明 \(\mathbf{K}\) 的代数闭包在 \(\mathbf{K}\) 上是无限扩张。
  \end{enumerate}
\end{enumerate}

\begin{enumerate}
\def\labelenumi{(\alph{enumi})}
\setcounter{enumi}{1}
\tightlist
\item
  设 K 为域  的有限生成扩张。证明若 K 在 \(K_{0}\) 上不代数，则 K 上存在离散赋值 v，使得在 \(K_{0}\) 上 。\(v(x) = 0\)\(K_{0}\)。
\end{enumerate}

